% Source of truth: preprint6_en.md (clean, world-facing). This file was GENERATED
% from it for v1.0.0-v1.0.1; the generating wrapper (build.py, in the generating
% session's scratchpad, after the design of active-v10.2/src/tools/md_to_tex.py)
% is LOST, and no surviving generator reproduces this layout.
%
% v1.0.2 therefore carries HAND-MIRRORED edits: the two corrections of that
% release (erratum E3 -- the scope p not= 1 (mod 8) of the section law, in five
% places; clarification C2 -- Lemma 6(b) "primitive") were applied here by hand
% to match preprint6_en.md exactly, and the pair was verified by grep over the
% list of replacements. See history/ERRATA.md.
%
% Do not hand-edit further without the same verification; restoring the build
% wrapper is the standing fix.
\documentclass[11pt,a4paper,onecolumn]{article}
\usepackage[T1]{fontenc}
\usepackage[utf8]{inputenc}
\usepackage{lmodern}
\usepackage{amsmath,amssymb,amsthm}
\usepackage[margin=2.4cm]{geometry}
\usepackage{microtype}
\usepackage{longtable,booktabs}
\usepackage[hidelinks]{hyperref}
\setlength{\parskip}{0.55em}
\setlength{\parindent}{0pt}
% The book-emoji (source marker for "imported/classical, cited not measured")
% is not a LaTeX glyph: map it to a plain bracketed marker so it neither
% crashes the build nor silently disappears.
\newcommand{\bookmark}{\textbf{[cited]}}
\title{Crystallography from a lift condition\thanks{Geometry of the vacuum, floor 0}}
\author{Vladimir Sobol\thanks{ORCID~\href{https://orcid.org/0009-0006-9829-7931}{0009-0006-9829-7931}. Independent Researcher, Ukraine. Corresponding: \href{mailto:amperaggio@gmail.com}{amperaggio@gmail.com}.}}
\date{2026-08-23}
\begin{document}
\maketitle

\begin{abstract}
\emph{Finite point groups from a Heisenberg lift: a Lorentzian daughter metric, its positivity compactifier, and an integrality lattice generate crystallographic classes as a function of the observer, without a hand-chosen direction.}

Starting from a single deep wedge in the orthogonal Lie algebra so(5,3), the induced metric on the wedge's transverse quotient has signature exactly (3,1), and the equivariance of the degree-0 floor built over it \emph{forces} that same metric as the unique so(W)-equivariant bracket of family I. On this structure two compactifiers are measured --- a positivity compactifier cutting the Levi to a compact SO(3)$\times $SO(2) and an arithmetic thickness of the integrality stabiliser $\Gamma $ (measured as a thick/thin \emph{contrast}; on the split stratum the covolume of the arithmetic group O$^{+}$(I$_{3,1}$) = [3,4,4] is \emph{computed}, V = 0.0763304662, matching the orbit density 1/(48V) = 0.27294 against a direct count 0.2727 --- the general finite-covolume theorem remaining \bookmark{} Borel--Harish-Chandra) --- whose weave gives a finite-volume moduli. A canonical lattice does \emph{not} exist: the daughter's symmetry group acts with 3-dimensional orbits on an irreducible W, and finite states never yield one. What does exist is a \emph{class} of lattices (the lift condition), and --- shown on the standard $\mathbb{Z}$$^{4}$$\otimes $$\mathbb{Z}$$^{2}$ --- finite point groups arise as (compact stabiliser) $\cap $ (discrete arithmetic) per stratum, so the crystallographic class is a \emph{pair} ($\Lambda $, t) --- the same integral Lorentzian lattice I$_{3,1}$, read over different timelike sections t$^{\perp }$, realises all seven 3D holohedries (cubic, tetragonal, hexagonal, orthorhombic, monoclinic, trigonal (rhombohedral), triclinic), each identified by composition. Which holohedry a section carries is a \emph{classification}, not a representative-dependent lower bound: for p $\not\equiv $ 1 (mod 8), a rank-p section of I$_{p,1}$ exists iff its discriminant form is cyclic $\langle $1/n$\rangle $ (a quadratic-residue condition) --- the law's forward leg derived on every section, its reverse derived for p $\not\equiv $ 1 (mod 8); at p $\equiv $ 1 (mod 8) the reverse leg \emph{fails}, since the discriminant form does not fix the parity of the ambient. Read as a function of the parent, the map singles out the signature with exactly one timelike direction --- q$_{d}$ = 1, as geometry --- by a qualitative break (a definite section exists iff q$_{d}$ = 1), the unique signature with one observer to one finite crystal; the dimension d itself is \emph{not} singled out and remains an input. An author pre-registration that this variety would degenerate with dimension is a measured miss --- over the swept range the validated holohedry-image count [3, 7, $\geq $13] grows (a bounded result, not an asymptotic refutation). \textbf{No new constant is introduced; the single dimensional input is the ambient $\Lambda $-scale inherited by the programme.} \textbf{No physical reading of any object --- ``time'', ``energy'', ``mass'' --- is made or implied; where a compactness or covolume statement is classical it is marked \bookmark{}, not measured.}
\end{abstract}





\section*{\S{}0. Setup, provenance, and the reading discipline}

Let (V, $\eta $) be the real quadratic space $\mathbb{R}$$^{5,3}$, $\eta $ = diag(+1$\times $5, $-$1$\times $3). This ambient signature is the root of the assembly and is an \emph{input}: (5,3) is the unique small signature whose ``daughter'' (below) has signature (3,1) under the map (p,q) $\mapsto $ (p$-$2, q$-$2), and the map itself is re-derived in the folder, not assumed. The classical frames invoked throughout --- Levi--Malcev decomposition, sl(2)-module theory, Witt's theorem, Borel--Harish-Chandra finiteness of covolume, the crystallographic point-group orders --- are cited, not re-derived (\bookmark{}, multiplicity 0). This preprint lives on the \emph{daughter} and the \emph{degree-0 floor} --- the Jacobi bridge of \S{}2 --- and is floor 0 of the series; the earlier floor $-$1 work of the series (the dial factory) is a different level and is not carried here.

\textbf{Terminology.} Throughout, ``lattice'' means a discrete subgroup of translations (a Bravais lattice). The programme's earlier ``four-leaf'' is a frame primitive --- three axes spread in one volume and one orthogonal timelike direction, whose symmetry is a point group and whose periodicity is a consequence, not a datum --- a different object, not addressed here (\bookmark{} rhyme).

\textbf{Positioning.} Within our programme this is a side line, not the main thread --- the programme's central aim lies elsewhere. We present it in its own right because it turned out to be substantive: a self-contained arithmetic account of the (3,1) daughter's crystallography --- timelike sections of an integral Lorentzian lattice read through their discriminant forms, the crystallographic class as a pair ($\Lambda $, t), finite point groups as (compact positivity stabiliser) $\cap $ (arithmetic group) --- with no new constant introduced. It sits within a long and deep line of work on integral lattices, discriminant forms, and crystallographic groups --- Nikulin, Conway--Sloane, Vinberg, Brown--B\``{u}low--Neub\''{u}ser--Wondratschek--Zassenhaus among others --- on which it draws and which it cites throughout; what we add is a particular vantage and a few honestly bounded observations, not a new foundation. Its bridge to the programme's physical thesis is left open and stated as debt in \S{}7.

\textbf{Scope boundary (local--global).} The law of \S{}9 (Theorem 1) places this paper squarely inside a local-to-global problematic: the discriminant form is glue data, and the obstruction removing the hexagonal class at rank 2 is a local (quadratic-residue) condition deciding global realizability of a section. We state this and deliberately stop at the mathematical side of the divide. The reading of the same law as a statement about observers --- a local symmetry of a section versus the global realizability of a lattice containing it --- belongs to the physical interpretation of the programme, is not claimed here, and is left as a separate line of work (\S{}7). The boundary is drawn by design, not by ignorance of the question.


\section*{\S{}1. The daughter from the parent}
The parent space $\mathbb{R}$$^{5,3}$, $\eta $ = diag(+1$\times $5, $-$1$\times $3), enters as the root of the assembly: the signature (5,3) is forced by the target law (p,q) $\mapsto $ (p$-$2,q$-$2) --- it is the unique small signature whose image under that map is (3,1) --- and the map itself is re-derived here as an orthogonal decomposition $\eta $ $\cong $ H $\perp $ H $\perp $ $\eta $|W rather than carried in as a postulate. Scope: this covers the assembly as stated; the $\Lambda $-scale is inherited from the programme as an input, and no new constants are introduced.

A deep wedge X = u$\wedge $v on a totally isotropic 2-plane U = span(u,v) exists exactly when the Witt index of the ambient form is $\geq $ 2: at (5,3) the index is 3 and the wedge is present, while at (3,1) the index is 1 and it is absent --- there $\eta $(u,u) = $\eta $(v,v) = $\eta $(u,v) = 0 forces the residual form on u$^{\perp }$/$\langle $u$\rangle $ into the definite signature (2,0), leaving no second isotropic direction. Scope: the measurement is stated for $\mathbb{R}$$^{p,q}$ with p,q $\geq $ 2; at q = 3 only the single wedge is probed --- a saturated double-wedge would need p $\geq $ 4 $\wedge $ q $\geq $ 4 and is not measured here, so the fuller earlier claim stays outside this pass.

The daughter (3,1) IS: with W:= ker X / im X = U$^{\perp }$/U, the induced form $\eta $|W is nondegenerate, has signature exactly (3,1), and its radical is measured to coincide with U itself. This signature is constant across 24 group translates g$\cdot $X$\cdot $g$^{-1}$ with g = exp(so(5,3)) and across 24 further translates probing scale-blindness of the reading, but that constancy is measured on the exp-orbit --- the identity component of O(5,3) --- only; completeness of the orbit, i.e. that every deep wedge is reached this way, is Witt's theorem (\bookmark{}, multiplicity 0), cited here and not re-derived. Signature-independence from the choice of isotropic plane rests on the symbolic decomposition $\eta $ $\cong $ H $\perp $ H $\perp $ $\eta $|W, named in advance, before any number was computed; the behaviour at p or q > 5 is not measured (this caveat is on the daughter's re-derivation from the parent --- the signature sweep of \S{}6 separately measures the daughters (2,1), (3,1), (4,1) directly).

\section*{\S{}2. The degree-0 floor and the forcing of the metric}
The centralizer c of the wedge X in so(5,3) is fixed exactly: c = (sl(2) $\oplus $ so($\eta $|W)) $\ltimes $ heis(9), with the W-block of c lying in so($\eta $|W) to a measured residual of 5.2$\times $10$^{-16}$, at rank 18 with no remainder, and with two independent definitions of the Heisenberg subspace agreeing at rank 9 for their union. The scope carried with this measurement is the inventory of c itself --- sl(2) acts tracelessly on A and vanishes on the centre --- while the named decomposition of the bracket carried by heis(9) is deferred to the gluing statement below.

On this floor, equivariance does not merely admit a daughter metric but forces one: dim Hom$_{so(W)}$($\Lambda $$^{2}$(W$\otimes $A), $\Lambda $$^{2}$A) = 1, and the unique equivariant map is numerically proportional to $\eta $|W, with a residual of order 10$^{-16}$. The scope riding alongside this claim is load-bearing: equivariance was demanded only under so($\eta $|W), not under the larger sl(2) and not under n itself, so the dimension-1 result is the stronger statement, obtained under the smaller group. The $-$0.5 proportionality factor left open earlier was later closed as a normalization of the solver's centre coordinate, equal to exactly $-$1 in the named a$\wedge $b coordinate.

The Jacobi identity on this floor is transparent by construction: the image of the bracket lands in the centre, so triple brackets vanish identically by construction, not as a measured outcome. What is measured is a contrast --- a non-central bracket evaluating to 8.44$\times $10$^{2}$ set against the identically-zero central one --- and it is this contrast, not the identity itself, that carries the claim. Its scope is the carrier of the floor, n$_{0}$ = (W$\otimes $A)$\oplus $Z, where the bracket always lands in Z; outside central brackets the result asserts nothing.

The gluing of the floor to the parent bracket is exact: the heis-part bracket of c equals family I, [w$\otimes $a, w$'$$\otimes $b], with $\beta $ = $-$$\eta $|W at a residual of 0.0, in the named wedge convention, fixed in advance and matched to the code. The resulting ideal is non-abelian, [n,n] $\neq $ 0. The qualifier carried in the same breath as this claim is not optional: the sign $-$1 is a record of the named wedge convention, and a convention-free orientation discriminant remains an open, narrowed address.

Under the full Levi factor sl(2) $\oplus $ so($\eta $|W) --- rather than so($\eta $|W) alone --- the floor collapses to exactly the daughter metric and nothing else: family II dies, 6 $\to $ 0, while the daughter survives, 1 $\to $ 1, and this holds in both parent signatures (5,3) and (4,4). The scope is measured on two parents at n = 8, and the result sharpens the forcing statement: under the parent's full symmetry, no part of the type-II channel survives inside the equivariant kernel.

As a lemma bearing on that count, the family of type-II objects is 6-dimensional, factoring as 2$\times $3; the factor of 2 is the two-dimensionality of the commutant span\{1, $\bigstar $\}, where $\bigstar $$^{2}$ = $\pm $1 decides only $\mathbb{C}$ versus $\mathbb{R}$$\oplus $$\mathbb{R}$ inside that commutant. Measured directly across (3,1), (2,2), and (4,0), the factor of 2 is the same in every case --- it is blind to signature. This lemma's scope excludes the daughter, which does not use it, and it does not survive the full Levi factor, where the type-II family itself dies to 0; the sharpening from a signature-dependent to a signature-blind reading is recorded as a self-correction of the earlier formulation.

\section*{\S{}3. Why there is no canonical lattice}
No canonical lattice exists on the daughter space W, and this is one fact together with the metric-forcing daughter result carried by the same evidence, not two separate claims: the so($\eta $|W)-orbits are three-dimensional everywhere it was checked --- on timelike, spacelike, and lightlike vectors, and on 200 random probes, with orbit dimension measured minimum 3 and maximum 3 --- and W is irreducible under so($\eta $|W) (dim End(W) = 1, zero invariant lines); consequently any so($\eta $|W)-invariant set containing a nonzero vector must contain a three-dimensional manifold and cannot be discrete. Whatever forces the metric on this floor is, by the same orbit count, exactly what forbids a lattice on it. The scope of this result is the canonical, i.e. fully invariant, lattice: it says nothing about a lattice realized as a symmetry-breaking state rather than an invariant one, and the minimality of its three counted inputs --- scale, a splitting choice, and the lattice datum itself --- taken together as a set has not been established and stands as an open address. This is a measured absence, not an unfinished search: the mechanism is the collision of W's irreducibility with a three-dimensional orbit floor holding on every causal type probed, so the negative is a structural property of the invariant floor rather than a gap in coverage.

Passing to symmetry-breaking states reopens the question without adding a new input. Two generic elements of the translation floor W$\otimes $A were measured to have stabiliser dimension zero in the full Levi factor, uniformly across 50 probed configurations (range 0..0), so a class of states with a discrete stabiliser exists without a new handle. The scope of this measurement is four named families of configurations and three carrier types, not all configurations, and the qualifier that ``stabiliser dimension zero'' is a generic property rather than a lattice was carved before the numbers were taken, not added afterward. A narrowing recorded against this same node applies here: the zero measured is relative to the Levi factor only --- in the full centralizer the stabiliser of the winning class is seven-dimensional, so the discrete-looking result is a partial reading, not a lattice by itself.

Restricting to finite states forecloses the lattice for the winning class regardless of family. The point group of a generic member of that class has order 2, and this order-2 group is exactly the global kernel of the action on W$\otimes $A: by Schur's lemma (\bookmark{}, multiplicity 0), W and A both irreducible forces the kernel to \{$\pm $(I, I$^{-1}$)\}, so the effective point group is trivial. This is a completed enumeration, not a partial scan: 200 of 200 cases were covered, and completeness is established as a theorem --- a non-scalar matrix is non-derogatory, which forces its centralizer to be span\{I, T\} --- rather than by exhaustive search alone. The scope is finite states of the translation floor only; it is not a verdict on infinite states, which remain the one place a discrete translation group could still live.

\section*{\S{}4. A class of lattices from the bracket}
Two questions about the bracket on W$\otimes $A constrain any candidate lattice: whether a lattice on the vector space lifts to a discrete subgroup of the group generated by the bracket, and whether a canonical (functional-free) choice of such a lattice exists. Both are measured, not assumed.

A lattice $\Lambda $ $\subset $ W$\otimes $A lifts to a discrete subgroup of the group of the degree-0 floor if and only if M$^{\mathsf{T}}$BM is integral up to one common factor; this condition is real, not vacuous, and restrictive --- of 50 generic lattices tested, 0 pass. The bracket imposes on $\Lambda $ a condition that the vector-space reading alone does not see --- the object that this line of inquiry was set to confirm --- and the scope is bound to that reading: a classification is \emph{begun} --- the genus refinement of \S{}6/\S{}7 classifies the split section classes $\Lambda $$_{W}$ (at least 19 genera at |det| $\leq $ 8); the full $\beta $-integral family $\Lambda $ $\subset $ V remains the 26-dimensional type space of Theorem 1(3) --- and neither is completed. The centre step of the lift is not a new input: c(t$\Lambda $) = t$^{2}$$\cdot $c($\Lambda $) exactly (the generator of H$_{\Lambda }$ $\cap $ Z is $\tfrac{1}{2}$$\cdot $c($\Lambda $), Lemma 3), matching the action of the grading element $\delta $ already counted in the companion measurement. Two results are cited alongside this measurement purely as rhymes, at multiplicity 0, and are not leaned on for any step of the argument: the Stone--von-Neumann uniqueness class (\bookmark{}) and the classification of integer Lorentzian lattices (\bookmark{}).

Given that a lift condition exists, no canonical lattice can be selected from the bracket alone without a functional: the point group of a $\beta $-integral lattice is measured infinite, containing integer hyperbolic isometries --- maximal modulus |$\lambda $| = (2+$\sqrt{3}$)$^{2}$ = 7+4$\sqrt{3}$ $\approx $ 13.93 obtained by hand, with 3+2$\sqrt{2}$ a second independent witness. The implication carried here is a theorem, not a heuristic, and it carries the modulus explicitly: |$\lambda $| > 1 implies infinite order; only the converse --- infinite order implies |$\lambda $| > 1 --- is false, so a failed finite search proves nothing about order. The search window ($-$2, 2) used to exhibit the witness is a probe convention valid only for an existence claim, not for an exhaustive classification of the point group. Correspondingly, the composition of the point group inside the full Levi factor is closed on the split stratum --- O$^{+}$(I$_{3,1}$) = [3,4,4] (\S{}5/\S{}7) --- and off the split stratum stands as an open address.

A canonical discrete object is nonetheless present in the inventory, on a different carrier. Three independent legs --- scale-freeness at 5.6$\times $10$^{-16}$ against a moving contrast of 4.8$\times $10$^{1}$, closure at 0.0, and $\delta $-invariance --- measure $\mathbb{Z}$/4 = \{$\pm $1, $\pm $$\bigstar $\}. This object lives on $\Lambda $$^{2}$W and does not descend to W: the space of S $\in $ gl(W) sending each plane to its orthogonal complement is empty. A single side-rhyme is noted at multiplicity 0 and is not claimed as the same object: $\mathbb{Z}$/4 as the dual-circle number d = 3 of the prior assembly (\bookmark{}) sits on a different carrier entirely.

\section*{\S{}5. Two compactifiers}
Two independent measurements bear on the adapted family, and their splice is gated rather than assumed; the finite-volume conclusion itself is stated only at the weave below, under its two gating teeth.

\textbf{Positivity cuts the Levi to a compact stabiliser.} The stabiliser of a positive compatible complex structure J on (W$\otimes $A, $\eta $$\otimes $$\omega $), taken inside the full Levi so(3,1)$\oplus $sl(2), carries a Killing form that is negative definite ([$-$4..$-$4], dim = 4, type SO(3)$\times $SO(2)); an indefinite J$'$ instead retains a boost direction ([$-$4..+4], noncompact). The scope carried by this node is the resolution of the tension it raised: adaptedness to W$\otimes $A is a measured stabiliser type --- a stratum K --- not a selection and not a hidden input; the bare-$\sigma $ orbit space keeps a generic stratum with trivial stabiliser alongside it, and the two sit as rows of one type-map rather than one being crowned over the other.

\textbf{The arithmetic side is thick, and its covolume is now a number.} For $\Gamma $, the stabiliser of a $\beta $-integral lattice, the orbit density N(R)/vol(B$_{R}$) on H$^{3}$ does not decay (sitting in [0.26, 0.43] for R $\in $ [1.5, 4]), while both thin comparison subgroups tested fall off by a factor $\sim $ 7 over the same window --- a measured thick/thin contrast. Moreover all norm $-$1 vectors of I$_{3,1}$ form a single O$^{+}$-orbit on the future sheet (\textbf{Lemma 8 (one-orbit lemma; P6.Lem8) --- derived, one import (Hermite's constant)}\label{P6:Lem8}: t$\cdot $t = $-$1 $\Longrightarrow $ t$^{\perp }$ is a positive unimodular rank-3 lattice, hence $\cong $ $\mathbb{Z}$$^{3}$ --- by Hermite's bound its minimum norm is $\leq $ $\gamma $$_{3}$$\cdot $det$^{1/3}$ = 2$^{1/3}$ < 2, giving a norm-1 vector that splits off $\langle $1$\rangle $, iterated down ranks 2 and 1 ($\gamma $$_{2}$ = 2/$\sqrt{3}$, $\gamma $$_{1}$ = 1); the sole import is the value $\gamma $$_{3}$ = 2$^{1/3}$, \bookmark{} --- so I$_{3,1}$ $\cong $ $\mathbb{Z}$$^{3}$ $\oplus $ $\langle $$-$1$\rangle $, and any two such splittings are O(I$_{3,1}$)-conjugate), so N$_{arith}$ is the single O$^{+}$-orbit of e$_{4}$, with Stab$_{O+}$(e$_{4}$) = O(3) (order 48) and the density converges to 1/(|Stab(e$_{4}$)|$\cdot $V) = 1/(48$\cdot $V) = 0.27294, where V = 0.0763304662 is the covolume of the Coxeter reflection group [3,4,4]; the direct count gives 0.2708, 0.2736, 0.2727 at R = 5, 6, 7 --- measured 0.2727 vs predicted 0.27294. So the finite covolume, formerly only cited from Borel--Harish-Chandra (\bookmark{}), is here computed for the split lattice.

\textbf{The weave.} Under the adapted moduli G/K = H$^{3}$$\times $$\mathbb{H}$$^{2}$ together with a thick $\Gamma $, the quotient (H$^{3}$$\times $$\mathbb{H}$$^{2}$)/$\Gamma $ has finite volume, so a functional would have a finite-volume compact to live on --- without any functional being constructed. This splice is gated on two teeth (the positivity tooth and the thickness tooth) and its line prints only when both are green; no functional is introduced, and no direction is chosen by hand --- this is the scope carried by this splice itself. As with the arithmetic-thickness reading, the covolume-finiteness ingredient entering this splice remains \bookmark{} (Borel--Harish-Chandra, multiplicity 0) in general; on the split stratum it is computed (V above).

\section*{\S{}6. The type map: crystallographic class as a pair ($\Lambda $, t)}
The finite point groups of the standard lattice $\mathbb{Z}$$^{4}$$\otimes $$\mathbb{Z}$$^{2}$ in the assembly are not imported from crystallography but \emph{derived} as intersections: each is K$_{J}$ $\cap $ $\Gamma $ = (the compact positivity stabiliser of \S{}5) $\cap $ (the discrete arithmetic of \S{}5), taken per stratum of J, with no point chosen by hand. The strata carry dim K$_{J}$ $\in $ \{4, 2, 2, 0\}, and dimension alone does \emph{not} classify the type: the two dim-2 strata have different finite groups (|K $\cap $ $\Gamma $| = 32 versus 4), so the type is the conjugacy class of K$_{J}$, not its dimension. On the most symmetric (isotropic-W) stratum |K $\cap $ $\Gamma $| = 192 with effective order 96 = octahedral O $\cong $ S$_{4}$ (order 24) $\times $ $\mathbb{Z}$/4, while the generic stratum is trivial ($\pm $I only); an A-module datum splits the type further --- $\tau $ = i gives $\mathbb{Z}$/4, $\tau $ = $\rho $ gives $\mathbb{Z}$/6, generic gives $\mathbb{Z}$/2 --- and a hyperbolic integer isometry (spectral radius (2+$\sqrt{3}$)$^{2}$ = 7 + 4$\sqrt{3}$, the isometry of \S{}4) is excised, so the intersection is genuinely finite.

The scope of this map is load-bearing and must ride with it: it is the map of \emph{one} lattice --- the standard $\mathbb{Z}$$^{4}$ $\otimes $ $\mathbb{Z}$$^{2}$, laid in by hand --- read through the \emph{single} section e$_{4}$, so its ``cube only'' reading is an artefact of that one section, corrected below. Two boundaries carved as errata (E-3, E-4) travel in the title's own banner rather than being buried: the order-3 elements present in the isotropic group are the body-diagonal rotations of the octahedron, and the octahedral group O $\cong $ S$_{4}$ (order 24) and the tetrahedral T(12) share exactly those same eight order-3 elements, so order-3 is \emph{not} a tetrahedral marker --- the histogram (order-4: 88, order-12: 32) identifies the full octahedral group, and ``tetrahedron'' is refuted by that composition, not merely qualified. Moreover T$_{d}$ cannot be the maximal point group of any lattice at all: $-$I $\notin $ T$_{d}$ makes it non-centrosymmetric, hence not a holohedry (\bookmark{}, multiplicity 0), so the tetrahedral question is a \emph{state with a basis}, a different object, and is not measured here.

The crystallographic class is therefore a \emph{pair} ($\Lambda $, t): the same integral Lorentzian lattice I$_{3,1}$, read over different timelike sections t$^{\perp }$, realises all seven 3D holohedry types --- cubic, tetragonal, hexagonal, orthorhombic, monoclinic, trigonal, triclinic --- each identified by composition, with a smallest witness norm n$_{min}$ = |t$\cdot $t|: cubic 1, tetragonal 2, hexagonal 6 (witness t = (1,1,1,3), section A$_{2}$$\oplus $$\langle $2$\rangle $, |Aut| = 24), orthorhombic 7, monoclinic 11, trigonal 13, triclinic 43. (The vector (1,1,1,2) has t$\cdot $t = $-$1 and gives a \emph{cubic} section, |Aut| = 48 --- not hexagonal; the earlier reading is corrected.) The hexagonal holohedry gives the earlier ``hexagonal trivalent'' hypothesis a measured W-trace --- a rhyme, not a bridge. Which holohedries occur is now the classification of Theorem 1(7) --- for p $\not\equiv $ 1 (mod 8), and here p = 3, a section exists iff its discriminant form is a cyclic $\langle $1/n$\rangle $ --- not a box-truncated lower bound; and the smallest witness norms n$_{min}$ above are \emph{exact}, established by exhaustive search over det $\leq $ 43 (not a swept-box lower bound). The count of $\beta $-integral \emph{classes} is refined too: the genus of the class refines the earlier SNF proxy --- the eleven SNF buckets at |det| $\leq $ 8 split into at least 19 genera (distinct discriminant bilinear forms), hence at least 19 classes; the exact genus count is parked on a 2-adic parity subtlety (new address ADDR-genus-exact), and the odd-lattice parity note is reconciled at the lit-gate.

Read as a function of the parent (p, q), the map singles out one signature by a qualitative break rather than by any count. A timelike section is definite --- hence carries a finite holohedry --- if and only if the daughter has exactly one timelike direction (q$_{d}$ = 1, i.e. parent q = 3): a timelike t removes one negative direction, leaving t$^{\perp }$ of signature (p$-$2, q$_{d}$ $-$ 1), definite exactly when q$_{d}$ = 1. This is the unique signature where one observer corresponds to one finite crystal; q$_{d}$ = 0 needs no observer (the whole lattice is definite) and q$_{d}$ $\geq $ 2 needs a rational negative q$_{d}$-plane --- so the assembly, asked for the observer-to-crystal correspondence, points at q$_{d}$ = 1 as geometry, a first touch of that correspondence from inside rather than as a postulate. The boundary is stated in the same breath and is essential: this singles out q$_{d}$ = 1, \emph{not} d = 3 --- within the q$_{d}$ = 1 column the daughters (3,1) and (4,1) both have definite sections and both exhibit a hexagonal section (while (2,1) has none --- the quadratic-residue obstruction, the first negative of the chain), so d itself is not singled out among q$_{d}$ = 1, and the d = 3 part remains an honest import. Two guards keep the reading honest: an author pre-registration that the variety would ``degenerate to elliptic'' with rising dimension runs opposite to the corrected measurement over the range swept --- the validated holohedry-image count [3, 7, $\geq $13] grows --- so the pre-registered \emph{direction} fails where measured; larger dimension is not measured, so this is a bounded result, not an asymptotic refutation; and the positivity compactifier is signature-blind (its stabiliser is compact at every q$_{d}$ $\in $ \{1,2\}), so the singling-out of q$_{d}$ = 1 lives in the section geometry, not in compactness.

\section*{\S{}6a. Independence of the two conditions, and calibration on known cases}

The finite point groups of \S{}6 arise as an \emph{intersection} --- a compact positivity stabiliser with a discrete arithmetic group --- and that phrasing earns its weight only if both conditions are necessary. A single detector, \{U $\in $ PointGroup($\Lambda $): UJ = JU\}, applied to three worlds shows they are: dropping either condition breaks ``finite and nontrivial''. With a $\beta $-integral lattice but an \emph{indefinite} J$'$ the commuting set is infinite --- 204 commuting elements, among them P $\otimes $ I of infinite order --- so positivity is necessary; with a positive J$_{0}$ but a \emph{generic} (non-lifting) lattice only $\pm $I of 1188 candidates survive (effective order 1, trivial) --- so integrality is necessary; and with both the group is finite and nontrivial (192, effective 96, element orders [1,2,3,4,6,12] as in \S{}6). The three worlds are named cases on $\mathbb{Z}$$^{4}$ $\otimes $ $\mathbb{Z}$$^{2}$ and a generic M$\cdot $$\mathbb{Z}$$^{8}$, not a sweep of all combinations, and the value of the trichotomy is that one code runs all three --- it repackages two already-established facts (the finite 192 and the infinite hyperbolic order) into a single detector, rather than adding a new measurement.

Two known cases calibrate the construction, and the two carry \emph{different} epistemic weight --- a distinction drawn explicitly. The automorphism datum is a genuine blind reproduction: the A-module of \S{}6 was computed by integer enumeration of K $\cap $ $\Gamma $ --- a scan of the source for closed-form CM tokens is empty --- and it returns $\mathbb{Z}$/4, $\mathbb{Z}$/6, $\mathbb{Z}$/2 at $\tau $ = i, $\rho $, generic, which are exactly the automorphism groups Aut(E$_{\tau }$) of the corresponding elliptic curves (\bookmark{}, Silverman, \emph{Arithmetic of Elliptic Curves} III.10); neither source saw the other, so this is a reproduced known result. The polarisation datum is weaker and is stated as a \emph{calibration}, not a reproduction: the lift form was \emph{built} from prescribed elementary divisors, so recovering those divisors in pairs (d, d) is arithmetic of the Frobenius normal form, not an independent recovery of Mumford's theorem --- our lift datum (c, Smith(B/c)) coincides with Mumford's polarisation type, and the non-tautological content is only the identification of the two and the distinction of principal from non-principal ((1,1) versus (1,4)). The finite group orders on the \emph{adapted} split strata --- 192, 288, 96, matched by the full order-histogram --- are closed by Lemma 6: each factors as the tabulated cubic holohedry (order 48) times Aut(E$_{\tau }$) $\in $ \{$\mathbb{Z}$/4, $\mathbb{Z}$/6, $\mathbb{Z}$/2\} at $\tau $ = i, $\rho $, generic, none new, with Brown--B\``{u}low--Neub\''{u}ser--Wondratschek--Zassenhaus (1978, \bookmark{}) as the ceiling reference. The two dim-2 strata (orders 32 and 4) are $\tau $ = i and \emph{non-adapted} (their W-factor is not a rank-3 section holohedry --- the order-4 histogram reads 32 = D$_{4}$ $\otimes $ $\mathbb{Z}$/4, not tetragonal $\otimes $ $\mathbb{Z}$/2), so they lie outside Lemma 6; off the split stratum no such check is made, and no order is claimed new.

\section*{\S{}7. Open questions (the graph's addresses, verbatim --- not promises)}
These are the open addresses of the folder's graph, stated verbatim as questions, not as promises. Each is a named place where the assembly does \emph{not} yet reach; none is claimed as in progress unless a probe is reserved.

\textbf{The physics bridge is not built --- and this is deliberate.} No object here is read physically: nothing is called time, energy, mass, or a collective/ordered phase. Such readings would live not on the vector daughter W but on a \emph{spinor floor} --- so(3,1) carries a two-dimensional complex representation that the vector W does not see, and Dirac-type operators, a mass gap, and any ordered phase (for instance a superconducting one) would have to be built there. Whether the discrete structure of this preprint descends to that floor is a separate question, not attempted here and reserved by the author's word (\bookmark{} rhyme). We state this as an open debt so the reader is not left to infer a physical content the construction does not claim. The reading of \S{}9 (Theorem 1)'s law as a local-versus-global statement about observers is likewise a boundary drawn by design, not a gap, and is left to the programme's physical interpretation (see the scope boundary in \S{}0).

\textbf{On the type map and lattice classes.} The per-\emph{representative} set of holohedries reachable within a fixed box, \texttt{type\_set}, is still only a \emph{lower bound} --- it depends on the representative and on the box of sections; the covering radius of sections, and a \emph{proven} cross-class split, are not shown. (By contrast the n$_{min}$ table of \S{}6 is now \emph{exact} --- the smallest section norm realising each holohedry, by exhaustive search over det $\leq $ 43.) On the \emph{count of classes}: the SNF proxy is not the exact GL(4,$\mathbb{Z}$) classification --- it does not see parity (counter-pair diag(1,1,1,$-$3) $\perp $ A$_{2}$$\oplus $U, a \bookmark{} rhyme via the external skeleton); the genus of the $\beta $-integral class \emph{refines} it --- the eleven SNF buckets at |det| $\leq $ 8 split into at least 19 genera (distinct discriminant bilinear forms), hence at least 19 classes, with the |det| = 1 case closed as a single class. Whether that lower bound is an equality --- whether each genus is a single class and whether the exact genus count is 19 --- is parked on a 2-adic parity subtlety and reserved as \textbf{ADDR-genus-exact}. And a methodological rule adopted throughout: any criterion on which the word ``complete'' depends needs a negative control \emph{on the criterion itself} --- a deliberately narrowed criterion whose count must fall --- to be entered into every tact that carries an enumeration.

\textbf{On the point group and its composition.} On the \emph{split} stratum this is now closed: the $\beta $-integral lattice's arithmetic point group is O$^{+}$(I$_{3,1}$) = the Coxeter reflection group [3,4,4] (covolume 0.0763304662, \S{}5/Lemma 5), so its composition is identified there. Off the split stratum the composition inside the full Levi is not decomposed --- the infinitude is measured, not the make-up. The tetrahedral question is a \emph{state with a basis} (lattice + shift; e.g. diamond = fcc + a two-point basis) --- a different object than a lattice, not measured, and reserved by the author's word. Minor: the freeness of $\langle $P,P$'$$\rangle $ is unchecked, and in the window R$\leq $4 the free orbit coincides with the cyclic one, so exponential growth is not exhibited --- the load-bearing contrast does not need it.

\textbf{On the floor and its module.} The decomposition of the type-II family (F into three $\eta $$'$ / the c--$\bigstar $c pairs by name) is not done --- only the two-dimensionality of the commutant span\{1,$\bigstar $\} is measured, blind to signature. Equivariance under n itself (the adjoint action of the Heisenberg part) was not required --- the full centralizer c did not enter the equivariance. A convention-free discriminant of the daughter's \emph{orientation} --- the sign is a record of the named wedge convention, the definiteness is a fact --- is open, if it exists at all; a \bookmark{} rhyme to the assembly's sign-non-derivability is noted, not claimed.

\textbf{On the carrier and the range.} Other carriers of the daughter (not ker/im) --- whether they give the same --- were not asked. Signatures p or q > 5 are not measured; the saturated double-wedge does not exist at q = 3, so the full subject of that earlier claim is not measured in the assembly. The non-generic layer of the state class is not described, and pairs of more than two elements were not measured separately.

\textbf{On the inputs and what lies beside them.} The minimality of the three inputs (scale $\cdot $ frame $\cdot $ lattice) \emph{as a set} is not proven --- a smaller set may buy the same. Combinations of the inventory's slots were not swept --- each slot was asked in isolation. The spinor floor named in the physics-bridge debt above (the two-dimensional complex representation of so(3,1) that the vector W does not see) is the natural place to ask whether the discrete lives there --- not asked, pending a separate word of the author (ADDR-spinor).

\textbf{On the parent-dimension disposition.} The mark that would single out d = 3 (or (3,1) over (4,1), (5,1)) is absent among lattices --- the q$_{d}$ = 1 column is homogeneous --- so the candidate for such a mark lies beyond lattices, in a spinor floor (which motivates ADDR-spinor a second time). At q$_{d}$ $\geq $ 2 the observer is a rational negative q$_{d}$-plane in place of t: by Lemma 4b the adapted orbit is Gr$^{-}$$_{q_{d}}$(W) $\times $ $\mathbb{H}$$^{2}$ and by Lemma 5 $\Gamma $ is arithmetic there too, so the items (b)/(c) --- an adapted single orbit and a thick $\Gamma $ at q$_{d}$ $\geq $ 2 --- follow from those lemmas rather than being open; only the explicit finite-group count there is uncomputed (at q$_{d}$ = 1, compactness (a) was the part separately measured).

\textbf{On method (an instance of pre-registration).} One prediction of this line was registered before the numbers ran --- that the variety of holohedry types would degenerate toward the elliptic with rising dimension. Over the range swept the corrected measurement runs opposite: the validated holohedry-image count [3, 7, $\geq $13] grows. It is recorded not as a result but as an instance of discipline --- a pre-registered expectation whose direction fails where measured is reported as a failure --- with the honest bound that larger dimension is not measured, so the asymptotic prediction is not refuted, only its direction over the swept range. The earlier report rested on a contaminated count [4, 8, 19]; that count is withdrawn (its enumeration returned a subset of the automorphism group, not the group).


\section*{\S{}8. The core exact sequence and its classical readings}

The objects of \S{}\S{}2--6 all sit on one short exact sequence --- the degree-0 floor as a Heisenberg extension of its translation space by its centre,

\begin{verbatim}
        0  -->  Z  -->  H  -->  V  -->  0
                |       |       |
                |       |       \-  V = W (x) A  (the translation floor)
                |       \---------  H  (the degree-0 floor, non-abelian: [n,n] != 0)
                \-----------------  Z  (the 1-dimensional centre, spanned by X)
\end{verbatim}

The four data attach to this one sequence as, respectively --- Malcev: rational bracket $\to $ lattice existence $\cdot $ SvN: central character $\to $ representation $\cdot $ Mumford: alternating form $\to $ polarisation type $\cdot $ ours: integrality + positivity $\to $ K$\cap $$\Gamma $.

Several classical structures are \emph{readings} of this one core --- each a different datum attached to the same sequence, none a functor beyond what is stated. We tabulate the readings and their outputs; the classical rows are imported (\bookmark{}, multiplicity 0), the last row is the present construction.

\begingroup\small
\begin{longtable}{p{1.7cm}p{2.5cm}p{3.3cm}p{3.1cm}p{3.2cm}}
\toprule
\textbf{reading} & \textbf{classical datum (\bookmark{})} & \textbf{our datum} & \textbf{output} & \textbf{checked} \\
\midrule\endhead
\textbf{Malcev 1949} & rational structure constants of the nilpotent Lie algebra & the bracket $\beta $ = $-$$\eta $|W is rational (indeed integral in the named convention) & a lattice exists in H (\bookmark{} Malcev) & cited (\bookmark{}) \\
\textbf{Stone--von-Neumann} & the central character of H & the centre Z is 1-dimensional, spanned by X & the irreducible unitary representation is unique (\bookmark{}) & cited (\bookmark{}) \\
\textbf{Mumford 1966} & the elementary divisors of the symplectic form on V & the Smith invariant factors of the lift form B & the polarisation type = those divisors (\bookmark{} Mumford) & calibrated (B-1, \S{}6a) \\
\textbf{ours (arithmetic $\Gamma $, discrete)} & --- & the integrality condition M$^{\mathsf{T}}$BM integral up to one factor & $\Gamma $ = stabiliser of a $\beta $-integral lattice is discrete (continuous part 0) & measured \\
\textbf{ours (arithmetic $\Gamma $, thick)} & --- & the orbit-density thick/thin contrast on H$^{3}$ & $\Gamma $ is thick (orbit density does not decay); covolume of [3,4,4] & measured contrast; covolume computed V = 0.0763304662, density 1/(48V) = 0.27294 vs count 0.2727 (was \bookmark{} Borel--HC) \\
\textbf{ours (observer stabiliser K$_{t}$)} & --- & positivity + a timelike section t & a compact stabiliser K$_{t}$, and per stratum the finite group K$_{J}$ $\cap $ $\Gamma $ & A-module blind (B-2, \S{}6a); orders split-closed by Lemma 6 (\S{}6a) \\
\bottomrule
\end{longtable}
\endgroup

\textbf{Theorem-level bridges (one per neighbour, with identifier).} Each classical reading rests on a named theorem, cited, not re-derived (\bookmark{}, multiplicity 0):

\begin{itemize}
\item \textbf{Malcev 1949} (Izv. AN SSSR 13): a simply-connected nilpotent Lie group admits a lattice if and only if its Lie algebra has rational structure constants --- the theorem under the ``our datum'' of row 1; here the bracket is integral in the named wedge convention, so the hypothesis is met.
\item \textbf{Stone--von-Neumann}: the irreducible unitary representation of the Heisenberg group is unique given the central character --- the uniqueness under row 2; our centre is measured 1-dimensional.
\item \textbf{Mumford 1966} (Invent. Math. 1, 287--354; Tata Lectures on Theta): the type of a polarisation is the tuple of elementary divisors of the alternating form --- the statement under row 3, reproduced by blind computation of the Smith invariant factors, which recovered (1),(2),(3),(1,1),(1,2),(2,2),(1,4).
\item \textbf{Borel--Harish-Chandra 1962} (Ann. Math. 75, 485--535): an arithmetic subgroup of a semisimple group has finite covolume --- the theorem granting the finite volume used in \S{}5, imported and not measured here (on the split stratum the covolume is computed, \S{}5).
\end{itemize}

\textbf{Structural analogues (weaker than a bridge --- no shared theorem is claimed).} Two recent lines share \emph{structure} with the present construction without a common theorem: the code-CFT construction of Dymarsky and Shapere (JHEP 03, 2021, 160) pairs integrality of a Lorentzian lattice with a positive-definite splitting to select points, which is structurally our intersection K $\cap $ $\Gamma $ but with a different output (CFT points, not crystallographic classes); and hyperbolic band theory (Maciejko--Rayan, Sci. Adv. 7, 2021; PNAS 119, 2022) takes the quotient H$^{2}$/$\Gamma $ of a Fuchsian group as a moduli object, of which the finite-volume moduli (H$^{3}$ $\times $ $\mathbb{H}$$^{2}$)/$\Gamma $ of \S{}5 is a three-dimensional, arithmetic relative. These are named as analogues, at multiplicity 0, and no result of \S{}\S{}2--6 depends on them.

\textbf{Relation to prior work.} Symmetry classifications of full spacetime lattices go back to Janner and Ascher's relativistic crystallography (Physica 45, 1969) and continue in the spacetime-crystal literature (Xu--Wu 2018; Li et al. 2020; Gopalan 2021); those works classify the symmetry group of one spacetime pattern, with the time axis fixed. Hyperbolic band theory (Maciejko--Rayan 2021, 2022; Boettcher et al. 2022) uses the quotient H$^{2}$/$\Gamma $ of a Fuchsian group as a moduli object; the finite-volume moduli (H$^{3}$$\times $$\mathbb{H}$$^{2}$)/$\Gamma $ of \S{}5 is its three-dimensional, arithmetic relative. In code CFTs (Dymarsky--Shapere 2021) integrality of a Lorentzian lattice is paired with a positive-definite splitting to select points; and the type of a polarisation in Mumford's theta-group theory (1966) is the closest mathematical ancestor of a finite datum arising from Heisenberg integrality together with positivity. The mechanism of cutting a definite lattice as the orthogonal complement of a timelike vector in an integral Lorentzian lattice is itself classical --- it is Vinberg's guiding-vector construction (1972) and, in the isotropic form, Conway's w$^{\perp }$/w passage in II$_{25,1}$ (Conway--Sloane, ch. 26--27) --- and the finite groups it yields are the tabulated crystallographic ones; what is new here is of formulation, not of mechanism. To our knowledge, three such formulations do not appear in that literature: reading the crystallographic class of a fixed integral Lorentzian lattice as a \emph{function of the timelike section} (the class as a pair ($\Lambda $, t)); deriving the lattice class itself from the \emph{lift condition of a Heisenberg bracket}; and presenting the per-stratum finite point group as the intersection of a \emph{compact positivity stabiliser} with the arithmetic group. Each is a repackaging of the classical section mechanism, not a new group or a new theorem about lattices. For the split stratum this check is closed by Lemma 6: every point-group order factors as a tabulated 3D holohedry times Aut(E$_{\tau }$) --- none new, with Brown et al. (1978) as the ceiling reference; off the split stratum no check is made, and no claim of novelty is attached to the group orders themselves. The common thread across these readings is the Heisenberg core of \S{}8: Malcev rationality relates the bracket to lattice existence; the central character supplies the Stone--von Neumann representation-theoretic reading; Mumford's polarisation type reads the elementary divisors of the same alternating datum; and the present construction adds positivity and arithmetic intersection to obtain the finite point groups. These are distinct readings of the same central extension, not identifications of the resulting structures.


\section*{\S{}9. The formal spine --- definitions, lemmas, and the Main Theorem (Theorem 1)}

This section states formally what \S{}\S{}1--8 established, so the claim reads as a theorem with hypotheses. Each item carries its status: \textbf{derived} (a proof is given here, a few lines, needing no probe), \textbf{measured} (it rests on a probe), \textbf{\bookmark{}} (an imported theorem, multiplicity 0). Nothing is read physically. Component groups are named explicitly throughout --- O, SO, O$^{+}$, SO$^{+}$ --- and the central quotient factor $\pm $ is kept beside each group. Statements of this paper are cited across the series by their global identifiers P6.Def k, P6.Lem k, P6.Thm k.

\subsection*{Definitions}

\textbf{Definition 0 (the data).}\label{P6:Def0} Let (W, $\eta $) be the daughter: a real quadratic space of signature $\Sigma $ = (p$_{d}$, q$_{d}$), d := p$_{d}$ + q$_{d}$. Let (A, $\omega $) = ($\mathbb{R}$$^{2}$, $\omega $) be the sl(2)-module with its symplectic form. Put V := W $\otimes $ A and \textbf{B := $\eta $ $\otimes $ $\omega $}, an alternating nondegenerate form on V. The Heisenberg algebra is n := V $\oplus $ $\mathbb{R}$X with [v, v$'$] := B(v, v$'$)$\cdot $X; H := exp n is the simply-connected group (dim 2d+1), and Z := exp $\mathbb{R}$X its one-dimensional centre. The symmetry group is \textbf{G := O($\eta $) $\times $ SL(A)}, acting on V as g $\otimes $ h; it preserves B, and the action map G $\to $ Sp(V, B) has kernel $\langle $($-$1, $-$1)$\rangle $, so its image is G/$\langle $($-$1, $-$1)$\rangle $ $\subset $ Sp(V, B) (this $\pm $ is carried wherever G acts on V). Write $\mathcal{J}$ := \{J $\in $ End V : J$^{2}$ = $-$1, B(J$\cdot $, J$\cdot $) = B, B(x, Jx) > 0 for x $\neq $ 0\} for the Siegel space of compatible positive complex structures, and K$_{J}$ := Stab$_{G}$(J) $\subset $ G. Since the action G $\to $ Sp(V, B) has kernel $\langle $($-$1,$-$1)$\rangle $, the stabiliser subgroups live in G \emph{without} a quotient; we write $\bar{K}$$_{J}$ for the image of K$_{J}$ in Sp(V, B) --- a quotient of K$_{J}$ by the $\langle $($-$1,$-$1)$\rangle $ it contains --- and all /$\langle $($-$1,$-$1)$\rangle $ equalities below are stated for the barred (image) groups.

\textbf{Definition 1 ($\beta $-integrality --- the exact lift datum).}\label{P6:Def1} A lattice $\Lambda $ $\subset $ V (rank 2d) is \emph{$\beta $-integral} if B($\Lambda $, $\Lambda $) $\subset $ c$\mathbb{Z}$ for some c > 0; let c($\Lambda $) be the positive generator of B($\Lambda $, $\Lambda $) $\subset $ $\mathbb{R}$. Its \emph{type} is the tuple of Smith invariant factors (d$_{1}$ | \dots{} | d$_{d}$) of B/c($\Lambda $) on $\Lambda $ --- the polarisation type (\bookmark{} Mumford 1966).

\textbf{Definition 2 (adapted J; observer; section --- the classified stratum).}\label{P6:Def2} For a negative-definite q$_{d}$-plane N $\subset $ W and $\tau $ $\in $ $\mathbb{H}$ (the upper half-plane), set S$_{N}$ := (+1 on N$^{\perp \eta }$, $-$1 on N), so S$_{N}$ $\in $ O($\eta $) with S$_{N}$$^{2}$ = 1, and \textbf{J$_{N,\tau }$ := S$_{N}$ $\otimes $ J$_{\tau }$}. An \emph{observer} is a point N $\in $ Gr$^{-}$$_{q_{d}}$(W); for q$_{d}$ = 1 it is a timelike line $\mathbb{R}$t. Relative to a lattice $\Lambda $$_{W}$ $\subset $ W it is \emph{rational} when N is spanned by lattice vectors; the \emph{section} is \textbf{S($\Lambda $$_{W}$, N) := $\Lambda $$_{W}$ $\cap $ N$^{\perp \eta }$}, a positive-definite lattice of rank p$_{d}$, with \emph{holohedry} the full orthogonal group O(S($\Lambda $$_{W}$, N)). \textbf{The section classification (Theorem 1 below) is stated on the G-rational, split stratum $\Lambda $ = $\Lambda $$_{W}$ $\otimes $ $\Lambda $$_{A}$ only; the generic and non-split strata are not covered by it.}

\textbf{Definition 3 (G-rationality).}\label{P6:Def3} A lattice $\Lambda $ $\subset $ V is \emph{G-rational} if G is $\mathbb{Q}$-defined relative to $\Lambda $ --- i.e. $\Lambda $ carries a $\mathbb{Q}$-structure with respect to which G is defined over $\mathbb{Q}$. A \emph{split} $\Lambda $ = $\Lambda $$_{W}$ $\otimes $ $\Lambda $$_{A}$ (with $\eta $($\Lambda $$_{W}$, $\Lambda $$_{W}$) $\subset $ $\mathbb{Q}$) is one example, not the definition: non-split $\mathbb{Q}$-forms of G (through restriction of scalars) also occur. Put $\Gamma $$_{\Lambda }$ := Stab$_{G}$($\Lambda $) $\subset $ G (without a quotient --- for a split $\Lambda $, $\Gamma $$_{\Lambda }$ = O($\Lambda $$_{W}$) $\times $ SL($\Lambda $$_{A}$)), and write $\bar{\Gamma }$$_{\Lambda }$ for its image in Sp(V, B).

\subsection*{Lemmas}

\textbf{Lemma 1 (the metric is forced --- derived; blind to signature).}\label{P6:Lem1} On the degree-0 floor the daughter metric is not merely admitted by so($\eta $)-equivariance but forced, and this holds for any nondegenerate $\eta $ in any signature.

\emph{Proof.} Decompose $\Lambda $$^{2}$V = ($\Lambda $$^{2}$W $\otimes $ S$^{2}$A) $\oplus $ (S$^{2}$W $\otimes $ $\Lambda $$^{2}$A). Now so($\eta $) acts trivially on A, so the so($\eta $)-invariants of the target $\Lambda $$^{2}$A are all of $\Lambda $$^{2}$A = $\mathbb{R}$$\cdot $$\omega $. In the source, for d $\geq $ 3 the standard representation W of so($\eta $) is \emph{absolutely irreducible}, so End$_{so(\eta )}$(W) = $\mathbb{R}$ and the invariants of W $\otimes $ W $\cong $ End(W) are one-dimensional, spanned by $\eta $ $\in $ S$^{2}$W; this single word fixes both counts --- (S$^{2}$W)$^{so(\eta )}$ = $\mathbb{R}$$\cdot $$\eta $ and ($\Lambda $$^{2}$W)$^{so(\eta )}$ = 0. Hence Hom$_{so(\eta )}$($\Lambda $$^{2}$V, $\Lambda $$^{2}$A) = Hom$_{so(\eta )}$(S$^{2}$W $\otimes $ $\Lambda $$^{2}$A, $\Lambda $$^{2}$A) = $\mathbb{R}$$\cdot $($\eta $ $\otimes $ $\omega $), one-dimensional, and the unique so($\eta $)-equivariant bracket of family I is B = $\eta $ $\otimes $ $\omega $ up to scale. $\blacksquare $

The dimension count does not see the signature of $\eta $: nothing in the proof distinguishes (3,1) from (2,2) or (4,0), so the singling-out of (3,1) happens in the section geometry of Definition 2, not here. The numerical measurement dim Hom = 1 at residual 10$^{-16}$ is the shadow of this lemma; the bracket's $-$1 normalisation is fixed in the named wedge convention.

\textbf{Lemma 2 (no invariant discrete subgroup --- derived).}\label{P6:Lem2} For d $\geq $ 2, V = W $\otimes $ A admits no nonzero SO$^{+}$($\eta $)-invariant discrete subgroup; in particular no invariant lattice. A discrete translation group can appear only as a symmetry-broken state.

\emph{Proof.} A connected group fixing a discrete set fixes it pointwise, and fixing by the group SO$^{+}$($\eta $) is annihilation by its Lie algebra so($\eta $). So a nonzero SO$^{+}$($\eta $)-invariant discrete subgroup would contain a nonzero so($\eta $)-fixed vector v $\in $ V = W $\otimes $ A $\cong $ W$^{\oplus 2}$ (A trivial under so($\eta $)), hence a nonzero so($\eta $)-fixed vector in W. But W has no trivial so($\eta $)-subsummand: for d $\geq $ 3 it is absolutely irreducible (Lemma 1); for d = 2 the single generator of so($\eta $) --- so(2) at $\Sigma $ = (2,0)/(0,2), so(1,1) at $\Sigma $ = (1,1) --- acts on W with nonzero eigenvalues $\pm $1 (a rotation weight, or a boost eigenvalue) and so has no nonzero kernel, the case $\Sigma $ = (1,1) included. Either way (W)$^{so(\eta )}$ = 0 --- a contradiction. $\blacksquare $

The measured shadow: the so($\eta $)-orbit dimension is exactly 3 on every causal type probed and across 200 random probes, and W is irreducible (dim End(W) = 1), so an invariant set with a nonzero vector cannot be discrete. This is a structural absence, not an unfinished search --- the collision of W's irreducibility with a three-dimensional orbit floor.

\textbf{Lemma 3 (the lift is $\beta $-integrality --- derived).}\label{P6:Lem3} With c($\Lambda $) the positive generator of B($\Lambda $, $\Lambda $) (Definition 1), a lattice $\Lambda $ $\subset $ V generates a discrete subgroup (a lattice) H$_{\Lambda }$ := $\langle $exp $\Lambda $$\rangle $ of the Heisenberg group H if and only if $\Lambda $ is $\beta $-integral; then H$_{\Lambda }$ $\cap $ Z = exp($\tfrac{1}{2}$$\cdot $c($\Lambda $)$\cdot $$\mathbb{Z}$$\cdot $X), and c(t$\Lambda $) = t$^{2}$$\cdot $c($\Lambda $).

\emph{Proof.} H is 2-step nilpotent, so BCH truncates: exp u $\cdot $ exp v = exp(u + v + $\tfrac{1}{2}$B(u, v)$\cdot $X). For u, v $\in $ $\Lambda $ we have $-$u$-$v $\in $ $\Lambda $, and exp u $\cdot $ exp v $\cdot $ exp($-$u$-$v) = exp($\tfrac{1}{2}$B(u, v)$\cdot $X) $\in $ H$_{\Lambda }$; ranging over $\Lambda $, the words with zero V-part fill exactly exp($\tfrac{1}{2}$$\cdot $B($\Lambda $, $\Lambda $)$\cdot $X) = exp($\tfrac{1}{2}$$\cdot $c($\Lambda $)$\cdot $$\mathbb{Z}$$\cdot $X), so H$_{\Lambda }$ $\cap $ Z = exp($\tfrac{1}{2}$$\cdot $c($\Lambda $)$\cdot $$\mathbb{Z}$$\cdot $X) and H$_{\Lambda }$/(H$_{\Lambda }$ $\cap $ Z) $\cong $ $\Lambda $. The subgroup B($\Lambda $, $\Lambda $) $\subset $ $\mathbb{R}$ is finitely generated; it is discrete $\Longleftrightarrow $ cyclic $\Longleftrightarrow $ $\Lambda $ is $\beta $-integral, and otherwise dense in $\mathbb{R}$, whereupon H$_{\Lambda }$ $\cap $ Z is dense and H$_{\Lambda }$ is not discrete. So H$_{\Lambda }$ is discrete $\Longleftrightarrow $ $\Lambda $ is $\beta $-integral, and then it is a lattice in H. Finally B(tv, tv$'$) = t$^{2}$B(v, v$'$) gives c(t$\Lambda $) = t$^{2}$$\cdot $c($\Lambda $). $\blacksquare $

The measured shadow: of 50 generic lattices tested, 0 pass the integrality condition M$^{\mathsf{T}}$BM integral up to one common factor, and c(t$\Lambda $) = t$^{2}$$\cdot $c($\Lambda $) holds exactly, matching the grading element $\delta $. Malcev's theorem (\bookmark{} Malcev 1949, multiplicity 0) then supplies the lattice in H, its rational-structure-constants hypothesis met in the basis $\Lambda $ $\cup $ \{c($\Lambda $)X\}.

\textbf{Lemma 4ab (positivity gives a compact stabiliser --- derived).}\label{P6:Lem4} Write $\bar{K}$$_{J}$ for the image of K$_{J}$ = Stab$_{G}$(J) in Sp(V, B). (a) For every J $\in $ $\mathcal{J}$, $\bar{K}$$_{J}$ is compact --- and so is K$_{J}$, a finite $\langle $($-$1,$-$1)$\rangle $-extension of it --- for every J. (b) The adapted J$_{N,\tau }$ := S$_{N}$ $\otimes $ J$_{\tau }$ lies in $\mathcal{J}$, its image stabiliser is $\bar{K}$$_{J_{N,\tau }}$ = (O($\eta $|N$^{\perp \eta }$) $\times $ O($\eta $|N) $\times $ SO(2)$_{\tau }$)/$\langle $($-$1, $-$1)$\rangle $ of dimension p$_{d}$(p$_{d}$$-$1)/2 + q$_{d}$(q$_{d}$$-$1)/2 + 1, and its orbit is G$\cdot $J$_{N,\tau }$ $\cong $ G/K = Gr$^{-}$$_{q_{d}}$(W) $\times $ $\mathbb{H}$$^{2}$.

\emph{Proof.} (a) A J $\in $ $\mathcal{J}$ makes (V, J, B($\cdot $, J$\cdot $)) Hermitian, and the subgroup of GL(V) preserving both J and B is the compact unitary group U(V, J, B) $\cong $ U(d). The image of G in Sp(V, B) is closed: the image of a homomorphism of linear algebraic groups is a closed subgroup (standard, e.g. Borel, \emph{Linear Algebraic Groups}), and the image of G($\mathbb{R}$) is an open, finite-index subgroup of the real points of that image, hence closed. So $\bar{K}$$_{J}$ = (image of G) $\cap $ U(V, J, B) is closed in a compact group, hence compact; and K$_{J}$, a finite $\langle $($-$1,$-$1)$\rangle $-extension of $\bar{K}$$_{J}$, is compact with it. (b) Compatibility: S$_{N}$ $\in $ O($\eta $) and J$_{\tau }$ $\in $ Sp($\omega $) give J$_{N,\tau }$$^{2}$ = $-$1 with J$_{N,\tau }$ preserving B; positivity is the tensor $\eta $$_{N}$ $\otimes $ $\omega $$_{\tau }$ of two positive-definite forms, $\eta $$_{N}$ := $\eta $$\circ $S$_{N}$ being positive-definite since S$_{N}$ flips the sign on N. An element (g, h) $\in $ G stabilises S$_{N}$ $\otimes $ J$_{\tau }$ iff gS$_{Ng}$$^{-1}$ = $\lambda $S$_{N}$ and hJ$_{\tau }$h$^{-1}$ = $\lambda $$^{-1}$J$_{\tau }$ with $\lambda $$^{2}$ = 1 (as S$_{N}$$^{2}$ = 1). The value $\lambda $ = $-$1 is excluded on the A-side: an h with hJ$_{\tau }$h$^{-1}$ = $-$J$_{\tau }$ anticommutes with J$_{\tau }$, so has the form [[a, b], [b, $-$a]] with det = $-$(a$^{2}$+b$^{2}$) < 0, not in SL(2,$\mathbb{R}$). Hence $\lambda $ = 1, so g commutes with S$_{N}$ --- preserving both eigenspaces, g $\in $ O($\eta $|N$^{\perp \eta }$) $\times $ O($\eta $|N) --- and h $\in $ centraliser(J$_{\tau }$) = SO(2)$_{\tau }$; modulo $\langle $($-$1,$-$1)$\rangle $ this is the stated $\bar{K}$. The A-side exclusion of $\lambda $ = $-$1 holds even at p$_{d}$ = q$_{d}$, where the W-side alone would permit the swap N $\leftrightarrow $ N$^{\perp \eta }$, so a general $\Sigma $ is covered. Finally O($\eta $)/(O($\eta $|N$^{\perp \eta }$) $\times $ O($\eta $|N)) = Gr$^{-}$$_{q_{d}}$(W) and SL(A)/SO(2) = $\mathbb{H}$$^{2}$. $\blacksquare $

For $\Sigma $ = (3,1) the dimension is 3 + 0 + 1 = 4 with identity component SO(3) $\times $ SO(2), and G/K = H$^{3}$ $\times $ $\mathbb{H}$$^{2}$ (the moduli of \S{}5). The measured shadow: the Killing form of the positive stabiliser is negative-definite ([$-$4..$-$4], dim 4, SO(3)$\times $SO(2)), while an indefinite J$'$ keeps a boost ([$-$4..+4], noncompact). Uniqueness of the maximal-dimension adapted orbit --- with the strata dim K$_{J}$ $\in $ \{4,2,2,0\} as its shadow --- is not derived here and stands as an open address.

\textbf{Lemma 5 (arithmeticity --- derived (a,b); (c) via \bookmark{} Borel--Harish-Chandra; (d) via \bookmark{} Borel density).}\label{P6:Lem5} Let $\Lambda $ $\subset $ V be G-rational (Definition 3). Then $\Gamma $$_{\Lambda }$ is arithmetic in G, and the adapted moduli (G/K)/$\Gamma $$_{\Lambda }$ has finite volume; this holds only on the G-rational stratum.

\emph{Proof.} (a) By Lemma 1 the G-invariant alternating forms on V are one-dimensional, spanned by B; when G is $\mathbb{Q}$-defined relative to $\Lambda $ (Definition 3) that line is $\mathbb{Q}$-defined, so B = c$\cdot $B$_{\mathbb{Q}}$ for a rational form B$_{\mathbb{Q}}$, whence B($\Lambda $, $\Lambda $) $\subset $ c$\cdot $$\mathbb{Q}$ and $\Lambda $ is $\beta $-integral up to scale (a split $\Lambda $$_{W}$ $\otimes $ $\Lambda $$_{A}$ with $\eta $($\Lambda $$_{W}$, $\Lambda $$_{W}$) $\subset $ $\mathbb{Q}$ is one instance). (b) For the $\mathbb{Q}$-structure of $\Lambda $, $\Gamma $$_{\Lambda }$ = Stab$_{G}$($\Lambda $) = G $\cap $ GL($\Lambda $) = G($\mathbb{Z}$) is arithmetic in G by definition (in the split case G($\mathbb{Z}$) = O($\Lambda $$_{W}$) $\times $ SL($\Lambda $$_{A}$), an example). (c) [\bookmark{} Borel--Harish-Chandra 1962] an arithmetic subgroup of a semisimple group has finite covolume, so vol(G/$\Gamma $$_{\Lambda }$) < $\infty $, whence vol((G/K)/$\Gamma $$_{\Lambda }$) < $\infty $; the quotient is noncompact iff $\eta $$_{\mathbb{Q}}$ is isotropic (Godement's compactness criterion, \bookmark{}; for I$_{3,1}$ it is isotropic). (d) The finite-volume conclusion of (c) is moreover \emph{equivalent} to G-rationality. Here G = O(3,1) $\times $ SL(2,$\mathbb{R}$) is semisimple with no compact factors, so if $\Gamma $$_{\Lambda }$ has finite covolume then the Borel density theorem (\bookmark{}, applied to $\Gamma $$_{\Lambda }$ $\cap $ G$^{0}$, of finite index) makes $\Gamma $$_{\Lambda }$ Zariski-dense in G, forcing G to be $\mathbb{Q}$-defined relative to $\Lambda $ --- i.e. $\Lambda $ is G-rational (Definition 3); the converse is (b)--(c). So finite volume holds \emph{exactly} on the G-rational stratum, a countable, positive-codimension locus in the $\beta $-integral family; a generic $\beta $-integral $\Lambda $ --- carrying no $\mathbb{Q}$-structure --- has $\Gamma $$_{\Lambda }$ that is not a lattice in G (infinite covolume). Whether such a generic $\Gamma $$_{\Lambda }$ is moreover finite is an open address, whose measured shadow is the generic stabiliser-dimension 0 across 50 configurations. $\blacksquare $

For the split $\Lambda $ = I$_{3,1}$ $\otimes $ U the group $\Gamma $$_{W}$ = O(I$_{3,1}$), and O$^{+}$(I$_{3,1}$) is the Coxeter reflection group [3,4,4] whose covolume is 0.0763304662 --- no longer only cited from (c) but computed, so the BHC input becomes a number on this stratum. The measured thick/thin contrast of the orbit density sits against exactly this covolume.

\textbf{Lemma 6 (the point group is a product --- derived (a,b); measured factorisation).}\label{P6:Lem6} For a split lattice $\Lambda $ = $\Lambda $$_{W}$ $\otimes $ $\Lambda $$_{A}$ with $\Lambda $$_{W}$ $\eta $-integral: (a) $\Gamma $$_{\Lambda }$ = O($\Lambda $$_{W}$) $\times $ SL($\Lambda $$_{A}$), with image $\bar{\Gamma }$$_{\Lambda }$ = (O($\Lambda $$_{W}$) $\times $ SL($\Lambda $$_{A}$))/$\langle $($-$1, $-$1)$\rangle $ in Sp; (b) for an adapted J = S$_{t}$ $\otimes $ J$_{\tau }$ with \textbf{primitive} rational timelike t (the line $\mathbb{R}$t is the datum, so its generator may be taken primitive without loss of generality), the point group $\bar{K}$$_{J}$ $\cap $ $\bar{\Gamma }$$_{\Lambda }$ $\cong $ (O(L)$_{\pm }$ $\times $ Aut(E$_{\tau }$))/$\langle $($-$1, $-$1)$\rangle $, where L = $\Lambda $$_{W}$ $\cap $ t$^{\perp \eta }$ (a positive-definite section, Definition 2), O(L)$_{\pm }$ = \{$\varphi $ $\in $ O(L) : $\bar{\varphi }$ $\in $ \{$\pm $1\} on disc L\}, and Aut(E$_{\tau }$) = Stab$_{SL(\Lambda _{A})}$($\tau $) $\in $ \{$\mathbb{Z}$/4, $\mathbb{Z}$/6, $\mathbb{Z}$/2\}; moreover O(L)$_{\pm }$ = O(L) whenever $\varphi $(n) $\leq $ 2 (n = |t$\cdot $t| $\in $ \{1, 2, 3, 4, 6\}); in general [O(L) : O(L)$_{\pm }$] = |image of O(L) in ($\mathbb{Z}$/n)$^{\times }$/\{$\pm $1\}| --- e.g. triclinic n = 43 has O(L) = \{$\pm $1\}, so O(L)$_{\pm }$ = O(L) even though $\varphi $(43) = 42.

\emph{Proof.} (a) If g $\otimes $ h $\in $ GL($\Lambda $) with g $\in $ O($\eta $), h $\in $ SL(A), then after scaling by $\lambda $ both $\lambda $g and $\lambda $$^{-1}$h are integral, with det($\lambda $g)$^{2}$$\cdot $det($\lambda $$^{-1}$h)$^{4}$ = 1 forcing both unimodular and $\lambda $$^{2}$ = 1; hence $\Gamma $$_{\Lambda }$ = O($\Lambda $$_{W}$) $\times $ SL($\Lambda $$_{A}$) in G, its image $\bar{\Gamma }$$_{\Lambda }$ carrying the /$\langle $($-$1,$-$1)$\rangle $ (Definition 3). (b) An element (g, h) fixes J = S$_{t}$ $\otimes $ J$_{\tau }$ iff g $\in $ Stab$_{O(\Lambda _{W})}$($\pm $t) and h $\in $ Stab$_{SL(\Lambda _{A})}$($\tau $) = Aut(E$_{\tau }$). The gluing is explicit both ways: since $\Lambda $$_{W}$ is unimodular and t primitive with t$\cdot $t = $-$n, the finite-index sublattice L $\oplus $ $\mathbb{Z}$t $\subset $ $\Lambda $$_{W}$ realises $\Lambda $$_{W}$ as the graph of an anti-isometry $\gamma $ : disc(L) $\cong $ disc($\mathbb{Z}$t) = $\mathbb{Z}$/n. A g preserving the splitting acts as ($\varphi $ on L, $\varepsilon $ on $\mathbb{Z}$t), and g preserves $\Lambda $$_{W}$ $\Longleftrightarrow $ its induced maps ($\bar{\varphi }$, $\bar{\varepsilon }$) preserve $\gamma $, i.e. $\bar{\varphi }$ = $\bar{\varepsilon }$ under $\gamma $ --- this gives at once which pairs ($\varphi $, $\varepsilon $) extend to $\Lambda $$_{W}$ and the restriction of Stab$_{O(\Lambda _{W})}$($\pm $t) onto O(L). With $\varepsilon $ $\in $ \{$\pm $1\} on $\mathbb{Z}$t the image is O(L)$_{\pm }$ = \{$\varphi $ : $\bar{\varphi }$ $\in $ \{$\pm $1\}\}, and if ($\mathbb{Z}$/n)$^{\times }$ = \{$\pm $1\} (i.e. $\varphi $(n) $\leq $ 2) then O(L)$_{\pm }$ = O(L) automatically; in general O(L)$_{\pm }$ = O(L) $\Longleftrightarrow $ the image of O(L) in ($\mathbb{Z}$/n)$^{\times }$/\{$\pm $1\} is trivial --- so n = 43 with O(L) = \{$\pm $1\} qualifies despite $\varphi $(43) = 42. Passing to images gives $\bar{K}$$_{J}$ $\cap $ $\bar{\Gamma }$$_{\Lambda }$ = (O(L)$_{\pm }$ $\times $ Aut(E$_{\tau }$))/$\langle $($-$1, $-$1)$\rangle $. $\blacksquare $

The measured factorisation: |$\bar{K}$$_{J}$ $\cap $ $\bar{\Gamma }$$_{\Lambda }$| = |O(L)$_{\pm }$|$\cdot $|Aut(E$_{\tau }$)| reproduces the \textbf{three adapted strata} --- 192 = 48$\cdot $4 (cubic $\otimes $ $\mathbb{Z}$/4, $\tau $ = i), 288 = 48$\cdot $6 (cubic $\otimes $ $\mathbb{Z}$/6, $\tau $ = $\rho $), 96 = 48$\cdot $2 (cubic $\otimes $ $\mathbb{Z}$/2, generic), effective = pair/2 --- matched by the full order-histogram, not by |$\bar{K}$ $\cap $ $\bar{\Gamma }$| alone, and built as a Kronecker group P $\subset $ GL(6,$\mathbb{Z}$) acting on L $\otimes $ $\Lambda $$_{A}$ (the timelike line $\mathbb{R}$t is not tracked, so the sign $\varepsilon $ is read off separately). The two dim-2 strata (orders 32 and 4) are $\tau $ = i and \textbf{non-adapted} (dim K$_{J}$ = 2, and the W-factor is not a rank-3 section holohedry --- by the order-4 histogram, 32 = D$_{4}$ $\otimes $ $\mathbb{Z}$/4, not tetragonal $\otimes $ $\mathbb{Z}$/2, which would give 8 not 20 elements of order 4), so they lie outside the lemma. On the adapted split stratum this closes the tabulation question: every order is (a tabulated 3D holohedry) $\times $ Aut(E$_{\tau }$), none new.

\textbf{Lemma 7 (the reverse direction --- derived for p $\not\equiv $ 1 mod 8; \bookmark{} only at p $\equiv $ 1).}\label{P6:Lem7} Let L be a positive-definite lattice of rank p whose discriminant group is cyclic, disc(L) $\cong $ $\mathbb{Z}$/n, with discriminant bilinear form $\langle $1/n$\rangle $ --- equivalently, admitting an anti-isometry $\gamma $ : disc(L) $\cong $ disc($\langle $$-$n$\rangle $) (the quadratic-residue condition). If p $\not\equiv $ 1 (mod 8) then L is a primitive timelike section of I$_{p,1}$.

\emph{Proof.} The graph of $\gamma $ in disc(L) $\oplus $ disc($\langle $$-$n$\rangle $) is an isotropic subgroup, and the corresponding overlattice M with L $\oplus $ $\langle $$-$n$\rangle $ $\subseteq $ M $\subseteq $ (L $\oplus $ $\langle $$-$n$\rangle $)$^{\ast }$ is unimodular --- a direct check: |graph $\gamma $| = n and disc(L $\oplus $ $\langle $$-$n$\rangle $) = n$^{2}$, so disc(M) = n$^{2}$/n$^{2}$ = 1, no import needed. M has signature (p, 1), and L = ($\mathbb{Z}$t)$^{\perp }$ is primitive in it (an element (x, y) $\in $ M with y = 0 has $\bar{y}$ = 0, so $\bar{x}$ = 0 by injectivity of $\gamma $, whence x $\in $ L; symmetrically $\mathbb{Z}$t = M $\cap $ $\mathbb{Q}$t). An even indefinite unimodular lattice has signature $\equiv $ 0 (mod 8); here sig(M) = p $-$ 1, so if p $-$ 1 $\not\equiv $ 0 (mod 8) --- i.e. p $\not\equiv $ 1 (mod 8) --- M cannot be even, hence is odd. By the classification of indefinite \emph{odd} unimodular lattices (\bookmark{} Serre, \emph{A Course in Arithmetic}, ch. V; Milnor--Husemoller 1973 --- the one import), the odd indefinite unimodular lattice of signature (p, 1) is unique, M $\cong $ I$_{p,1}$ = $\langle $1$\rangle $$^{p}$ $\oplus $ $\langle $$-$1$\rangle $. Hence L is a primitive timelike section of I$_{p,1}$. $\blacksquare $

For the measured signatures p = 2, 3, 4, 5 (p $-$ 1 = 1, 2, 3, 4, none $\equiv $ 0 mod 8) the reverse direction is therefore \emph{derived}, not imported. At p $\equiv $ 1 (mod 8), i.e. p $\in $ \{1, 9, 17, \dots{}\}, M could be even and parity alone does not force M $\cong $ I$_{p,1}$ --- and there the reverse direction \textbf{fails}, not merely stays imported: E$_{8}$ $\oplus $ $\langle $2$\rangle $ (rank 9) has disc $\cong $ $\mathbb{Z}$/2 with $\bar{b}$ = 1/2 and so satisfies the condition, yet it is a section only of the even II$_{9,1}$ and not of I$_{p,1}$ at p = 9, while I$_{8}$ $\oplus $ $\langle $2$\rangle $ carries the same discriminant form and \emph{is} a section of I$_{9,1}$. The discriminant form alone does not fix the parity of the ambient.

\subsection*{Main Theorem (Theorem 1)}

\textbf{Main Theorem (Theorem 1) --- the classification law.}\label{P6:Thm1} Let (W, $\eta $) be the daughter of $\mathbb{R}$$^{5,3}$ ($\Sigma $ = (3,1)), V = W $\otimes $ A, B = $\eta $ $\otimes $ $\omega $, G = O(3,1) $\times $ SL(2,$\mathbb{R}$). Then:

\begin{enumerate}
\item \textbf{(forcing)} B is the unique so($\eta $)-equivariant bracket of family I up to scale, blind to signature [Lemma 1]; under the full Levi it is the only survivor.
\item \textbf{(no invariant lattice)} V carries no nonzero SO$^{+}$($\eta $)-invariant discrete subgroup; discreteness is possible only as a symmetry-broken state, and for finite states the effective point group is trivial [Lemma 2; Schur \bookmark{}].
\item \textbf{(lift)} $\Lambda $ generates a lattice in H $\Longleftrightarrow $ $\Lambda $ is $\beta $-integral; then H$_{\Lambda }$ $\cap $ Z = exp($\tfrac{1}{2}$c($\Lambda $)$\mathbb{Z}$X); the type is a G-invariant. Solving M$^{\mathsf{T}}$BM = c$\cdot $B$_{0}$ makes the $\beta $-integral lattices of a fixed type a single Sp(V, B)-orbit (stabiliser the arithmetic group of the type), so modulo G and scale the family has dimension dim Sp(2d) $-$ dim G $-$ 1 = d(2d+1) $-$ (d(d$-$1)/2 + 3) $-$ 1 = (3d$^{2}$+3d)/2 $-$ 4 (d = 4: 36 $-$ 9 $-$ 1 = 26) [Lemma 3].
\item \textbf{(compactifier)} every $\bar{K}$$_{J}$ is compact, and the adapted J = S$_{t}$ $\otimes $ J$_{\tau }$ form one orbit G/K = H$^{3}$ $\times $ $\mathbb{H}$$^{2}$ with K = SO(3) $\times $ SO(2) [Lemma 4ab; orbit-uniqueness not derived].
\item \textbf{(arithmeticity --- the scope of finite volume)} for G-rational $\Lambda $, $\Gamma $$_{\Lambda }$ is arithmetic and (H$^{3}$$\times $$\mathbb{H}$$^{2}$)/$\Gamma $$_{\Lambda }$ has finite volume; for split $\Lambda $ = I$_{3,1}$ $\otimes $ U, $\Gamma $$_{W}$ = O(I$_{3,1}$) and O$^{+}$(I$_{3,1}$) = [3,4,4] with covolume 0.0763304662; for $\beta $-integral but non-G-rational $\Lambda $, finite volume \textbf{does not occur} (infinite covolume, by Lemma 5(d)) [Lemma 5 + \bookmark{} Borel--Harish-Chandra].
\item \textbf{(finite point groups = a pair ($\Lambda $, t))} for every J $\in $ $\mathcal{J}$ and every $\Lambda $, the point group $\bar{K}$$_{J}$ $\cap $ $\bar{\Gamma }$$_{\Lambda }$ is finite (compact $\cap $ discrete in Sp); for split $\Lambda $ and adapted J with primitive timelike t, $\bar{K}$$_{J}$ $\cap $ $\bar{\Gamma }$$_{\Lambda }$ $\cong $ (O(L)$_{\pm }$ $\times $ Aut(E$_{\tau }$))/$\langle $($-$1,$-$1)$\rangle $ with L = $\Lambda $$_{W}$ $\cap $ t$^{\perp \eta }$ a positive rank-3 lattice, det L = |t$\cdot $t|. The crystallographic class is the pair ($\Lambda $$_{W}$, t); the holohedry is O(L). The minimal section norm realising each holohedry is \emph{exact}, n$_{min}$ = (cubic 1, tetragonal 2, hexagonal 6, orthorhombic 7, monoclinic 11, trigonal 13, triclinic 43), by exhaustive search over det $\leq $ 43 [Lemma 6].
\item \textbf{(the section law --- the heart)} for p $\not\equiv $ 1 (mod 8): L is realised by a timelike section of I$_{p,1}$ $\Longleftrightarrow $ disc L $\cong $ $\mathbb{Z}$/n is cyclic with split bilinear form $\langle $1/n$\rangle $ (b$_{L}$(g,g) $\equiv $ a/n, a a quadratic residue mod n); membership in the image is decided by the discriminant form (derived: the forward implication together with Lemma 7); the image itself --- [3,7,13] at p = 2,3,4 and the mechanism at p = 5 --- is measured. The law is stated through the discriminant \emph{bilinear} form b$_{L}$ (well-defined mod $\mathbb{Z}$ for any lattice, even or odd); the finer quadratic form q$_{L}$ (mod 2$\mathbb{Z}$) is not used, so the classification is parity-agnostic --- the odd case runs through Lemma 7's parity argument, not Nikulin's even-lattice theorems. The forward implication is derived --- w = v + t/n gives b$_{L}$(w,w) $\equiv $ 1/n, checked on every section; the reverse --- primitive embedding into the unimodular ambient --- is \emph{derived} for p $\not\equiv $ 1 (mod 8) by gluing plus the odd-unimodular classification (Lemma 7), covering all measured signatures p = 2,3,4,5; at p $\equiv $ 1 (mod 8) it \textbf{fails} --- witness E$_{8}$ $\oplus $ $\langle $2$\rangle $ versus I$_{8}$ $\oplus $ $\langle $2$\rangle $ at p = 9, same discriminant form, different ambient parity. The image [3,7,13] is recovered from the discriminant form alone at p = 2,3,4 (13/13, by line-identical reproduction), and the mechanism holds at p = 5 (267/267); the rank-2 hexagonal gap is the quadratic-residue obstruction 2 $\notin $ ($\mathbb{Z}$/3)$^{\times 2}$.
\item \textbf{(q$_{d}$ = 1)} a section by one rational vector is definite $\Longleftrightarrow $ q$_{d}$ = 1; for q$_{d}$ $\geq $ 2 the observer is a rational negative q$_{d}$-plane (a point of Gr$^{-}$, not a flag) [derived].
\end{enumerate}

\textbf{What Theorem 1 does not claim} (so the abstract cannot overreach): no physical reading of any object; not the finiteness of the $\beta $-integral class (it is 26-dimensional); not finite volume off the G-rational stratum; not the singling-out of d = 3 (only q$_{d}$ = 1, and even that as ``observer = a vector, not a plane''); not the exact count of $\beta $-integral classes (only at least 19 genera, hence at least 19 classes, at |det| $\leq $ 8 --- ADDR-genus-exact; the n$_{min}$ table itself is exact); and no new finite group (each reduces to a tabulated one).

$\bigstar $ \textbf{Reverse-direction: derived (Lemma 7), with one import.} The reverse leg of (7) --- that a rank-p lattice with cyclic $\langle $1/n$\rangle $ discriminant form is a section of I$_{p,1}$ --- is \emph{derived} for p $\not\equiv $ 1 (mod 8) in Lemma 7: the graph of the anti-isometry glues L $\oplus $ $\langle $$-$n$\rangle $ to a unimodular M of signature (p,1) (elementary, disc = n$^{2}$/n$^{2}$ = 1), and since sig(M) = p $-$ 1 $\not\equiv $ 0 (mod 8) forces M odd, the uniqueness of the indefinite odd unimodular lattice (\bookmark{} Serre, ch. V; Milnor--Husemoller --- the sole import) gives M $\cong $ I$_{p,1}$. This covers p = 2,3,4,5. At p $\equiv $ 1 (mod 8) --- where M could be even --- the leg does not merely stay \bookmark{}: it \textbf{fails}, witnessed at p = 9 by E$_{8}$ $\oplus $ $\langle $2$\rangle $ (disc $\cong $ $\mathbb{Z}$/2, $\bar{b}$ = 1/2, a section only of II$_{9,1}$) against I$_{8}$ $\oplus $ $\langle $2$\rangle $ (the same discriminant form, a section of I$_{9,1}$). Nikulin's embedding theorems (1.12.2, 1.14.4) are for \emph{even} lattices, so they are not what is used here; the odd case runs through the parity argument and Serre's odd classification.


\section*{Lit-gate (separate step, one pass, external-data protocol)}
Three questions, answered ``known $\perp $ ours``: (a) integer Lorentzian lattices and finite subgroups of O(3,1;$\mathbb{Z}$) --- what is classical; (b) ``crystal class as a function of the observer/section'' --- is it in the literature; (c) deriving a lattice class from a Heisenberg-lift condition --- is it in the literature. Until this gate ran, the phrase ``to our knowledge'' was forbidden anywhere in the text; the gate has now run, and the three bounded uses appear in Relation to prior work above --- no other novelty phrasing is admitted.


\emph{Series: Geometry of the vacuum, floor 0 (numbering fixed per preprint-1).}


\nocite{*}
\bibliographystyle{plain}
\bibliography{preprint6}
\end{document}
